\documentclass[11pt]{article}

\usepackage[T1]{fontenc}
\usepackage[utf8]{inputenc}
\usepackage[english]{babel}
\usepackage{amsmath,amssymb,amsthm}
\usepackage{booktabs}
\usepackage{geometry}
\usepackage{hyperref}
\usepackage{url}

\geometry{margin=1in}

\hypersetup{
  colorlinks=true,
  linkcolor=black,
  citecolor=blue,
  urlcolor=blue,
  pdftitle={Computational Extension of OEIS Sequence A226188},
  pdfauthor={Carlo Corti}
}

\newtheorem{theorem}{Theorem}
\newtheorem{proposition}[theorem]{Proposition}
\newtheorem{lemma}[theorem]{Lemma}
\newtheorem{corollary}[theorem]{Corollary}
\theoremstyle{definition}
\newtheorem{definition}[theorem]{Definition}
\theoremstyle{remark}
\newtheorem{remark}[theorem]{Remark}

\title{Computational Extension of OEIS Sequence A226188:\\
Certified Inverse Harmonic Thresholds up to \(n=67\)}

\author{Carlo Corti\\
Independent researcher\\
\texttt{carlo.corti@outlook.com}}

\date{June 26, 2026}

\begin{document}

\maketitle

\begin{abstract}
OEIS sequence A226188 is defined by
\[
  a(n)=\min\{k\geq 1: H_k>2n/3\},
\]
where \(H_k\) is the \(k\)-th harmonic number.  The local OEIS snapshot used for
this computation contained 27 known terms.  This paper reports a certified
extension through \(n=67\).  For every reported value \(k=a(n)\), the
certificate is the strict threshold condition
\[
  H_{k-1}\leq 2n/3 < H_k.
\]
The computation uses the classical asymptotic approximation only to locate a
small search neighborhood.  The final decisions are made by outward-rounded
MPFR/GMP interval evaluations, using direct interval summation for small
arguments and Euler--Maclaurin formulae for large arguments, followed by an
expanding bracket and binary search.  The resulting \(64\)-bit implementation
ends at
\[
  a(67)=14054172515346621401;
\]
the next index requires a wider integer type.  The new data also give a
finite numerical consistency check, through the computed range, for the OEIS
ratio behavior associated with the asymptotic limit
\(a(n+1)/a(n)\to \exp(2/3)\).
\end{abstract}

\noindent\textbf{2020 Mathematics Subject Classification.}
Primary 11Y55; Secondary 11B83, 11A25.

\noindent\textbf{Key words and phrases.}
OEIS, harmonic numbers, inverse harmonic thresholds, Euler--Maclaurin summation,
interval arithmetic, computational number theory.

\section{Introduction}

The harmonic numbers
\[
  H_k = \sum_{j=1}^k \frac{1}{j}
\]
form a strictly increasing sequence.  OEIS A226188 asks for the first index at
which this sum exceeds the arithmetic progression \(2n/3\):
\[
  a(n)=\min\{k\geq 1: H_k>2n/3\}.
\]
The entry previously listed the 27 values
\[
\begin{gathered}
1,2,4,8,16,31,60,116,227,441,859,1674,3260,6349,\\
12367,24088,46916,91380,177984,346666,675214,1315136,\\
2561536,4989191,9717617,18927334,36865412.
\end{gathered}
\]
It also records the asymptotic consequence
\[
  \frac{a(n+1)}{a(n)}\to \exp(2/3)=1.947734041054676\cdots,
\]
attributed in the local snapshot to a comment by Jean-François Alcover.
The same local snapshot places A226188 in a small family of harmonic-threshold
OEIS sequences, including A226187 for the threshold \(n/3\), A226189 for
\(\sqrt n\), and A226190 for \(\log n\) \cite{oeisThresholdFamily}.

The present computation extends the table by 40 further terms, from \(n=28\)
through \(n=67\).  It is a finite computational certification, not a new proof
of an infinite theorem about all later indices.  The asymptotic formula
\[
  a(n)\sim \exp(2n/3-\gamma)
\]
is used as a locator only.  No reported value is accepted because it is close
to the asymptotic estimate; each value is accepted only after the two
neighboring harmonic inequalities have been certified.

The main contributions are:
\begin{enumerate}
\item reproduction of the 27 OEIS terms as a required gate;
\item certified generation of \(a(n)\) for \(28\leq n\leq 67\);
\item a reproducible C17/MPFR/GMP implementation using an expanding bracket,
      binary search, and directed interval arithmetic;
\item an explicit \(64\)-bit boundary: the current implementation supports
      \(n\leq 67\), while \(n=68\) belongs to a wider-integer continuation.
\end{enumerate}

\section{Definition and Certification Predicate}

\begin{definition}
Let \(H_0=0\), and for \(k\geq 1\) let
\[
  H_k=\sum_{j=1}^k \frac{1}{j}.
\]
For \(n\geq 1\), let
\[
  a(n)=\min\{k\geq 1: H_k>2n/3\}.
\]
\end{definition}

Because \(H_k\) is strictly increasing in \(k\), the following equivalent
predicate is the natural certificate for a reported term.

\begin{proposition}
\label{prop:certificate}
For \(n\geq 1\) and \(k\geq 1\), \(k=a(n)\) if and only if
\[
  H_{k-1}\leq \frac{2n}{3}<H_k.
\]
\end{proposition}

\begin{proof}
If \(k=a(n)\), then the defining strict inequality gives \(H_k>2n/3\), while
minimality gives \(H_{k-1}\leq 2n/3\).  Conversely, if these two inequalities
hold, then for every \(j<k\) one has \(H_j\leq H_{k-1}\leq 2n/3\), so \(k\) is
the first index at which the harmonic sum exceeds the threshold.
\end{proof}

This strict/weak pair is important: the OEIS definition uses \(H_k>2n/3\), so
the lower neighbor must be allowed to equal the threshold.

\section{Analytic Evaluation Framework}

Direct summation to \(k\) is infeasible at the largest values in this paper.
The final term is about \(1.4\cdot 10^{19}\), so a certification method must
evaluate \(H_k\) without summing all reciprocals.

For large \(k\), the implementation uses the Euler--Maclaurin expansion and
the standard first-neglected-term remainder control for the logarithmic-gamma
and psi-function asymptotics; see DLMF Section~5.11, especially the
asymptotic material in Section~5.11(ii), in \cite{dlmf}.  The
Euler--Maclaurin intervals below are used only for \(k>10000\); smaller
arguments are handled by direct interval summation.  Here
\(\gamma=0.5772156649\cdots\) is the Euler--Mascheroni constant and \(\log\)
denotes the natural logarithm.  In the present integer setting
\(H_k=\psi(k+1)+\gamma\), where \(\psi\) is the digamma function, and the
interval with four Bernoulli correction terms used for production comparisons
is
\[
\begin{aligned}
H_k ={}& \log k+\gamma+\frac{1}{2k}
        -\frac{1}{12k^2}+\frac{1}{120k^4}
        -\frac{1}{252k^6}+\frac{1}{240k^8}+R_4(k),
\end{aligned}
\]
where \(R_4(k)\) denotes the remainder after the displayed correction terms and
\[
  |R_4(k)|\leq \frac{1}{132k^{10}}.
\]
This remainder estimate is the first-neglected-term bound from the cited DLMF
asymptotic expansion.  If \(S_4(k)\) denotes the displayed sum without
\(R_4(k)\), then EM4 is the outward-rounded interval
\[
  \left[S_4(k)-\frac{1}{132k^{10}},\,
        S_4(k)+\frac{1}{132k^{10}}\right].
\]
As a higher-order check, the program also evaluates the interval with six
Bernoulli correction terms
\[
\begin{aligned}
H_k ={}& \log k+\gamma+\frac{1}{2k}
        -\frac{1}{12k^2}+\frac{1}{120k^4}
        -\frac{1}{252k^6}+\frac{1}{240k^8}\\
       &-\frac{1}{132k^{10}}+\frac{691}{32760k^{12}}+R_6(k),
\end{aligned}
\]
where \(R_6(k)\) is the corresponding remainder after the displayed terms and
\[
  |R_6(k)|\leq \frac{1}{12k^{14}}.
\]
This bound is used in the same way: if \(S_6(k)\) denotes the displayed sum
without \(R_6(k)\), then EM6 is the outward-rounded interval
\[
  \left[S_6(k)-\frac{1}{12k^{14}},\,
        S_6(k)+\frac{1}{12k^{14}}\right].
\]

All endpoint operations are rounded outward with MPFR \cite{mpfr}, linked
against GMP \cite{gmp}.  Outward rounding means rounding lower endpoints toward
\(-\infty\) and upper endpoints toward \(+\infty\).  Thus each comparison
uses an interval \([L_k,U_k]\) known to contain \(H_k\).  Here
\([L_k,U_k]\) is the active direct-summation, EM4, or EM6 interval at the
current working precision.  The predicate
\(H_k>2n/3\) is certified true only when \(3L_k>2n\), and certified false only
when \(3U_k\leq 2n\).  Unresolved interval comparisons cause the program to
increase precision or fail closed.

\section{Sequence-Specific Structure and Cross-Checks}

A226188 is not a rare-object search.  Every positive \(n\) has exactly one
answer, and the difficulty is to identify the crossing point of a monotone
analytic function on a rapidly growing integer scale.

There is also a simple exact reduction on one residue class.  If
\[
  b(t)=\min\{k\geq 1:H_k>t\}
\]
is the inverse harmonic threshold at an integer argument \(t\), as in OEIS
A002387, then
\[
  a(3m)=b(2m) \qquad (m\geq 1),
\]
because \(2(3m)/3=2m\).  Thus the rows with
\(n=30,33,\ldots,66\) in the new range coincide with the integer-threshold
values \(b(20),b(22),\ldots,b(44)\), respectively.  They are therefore
external OEIS cross-check rows against A002387 rather than structurally new
fractional-threshold cases.  OEIS A004080 records the corresponding
weak-threshold integer variant \(H_k\geq t\), and the OEIS pages for these
integer-threshold sequences cite the classical computation of Boas and Wrench
on partial sums of the harmonic series \cite{oeisIntegerThresholds,boasWrench}.
The remaining 27 new rows, where \(n\not\equiv0\pmod 3\), are the genuinely
fractional-threshold cases for A226188 in this range.

The sequence-specific computational feature is that the asymptotic estimate is
extremely accurate near the relevant \(k\), but it is not itself a certificate.
In the completed run, the seed
\[
  k_0=\left\lfloor\exp(2n/3-\gamma)\right\rfloor
\]
was often equal or adjacent to the final value.  Nevertheless, the algorithm
does not use equality to the seed as an acceptance rule.  It expands a bracket
until it has a certified false endpoint and a certified true endpoint, and then
uses the monotonicity of \(H_k\) to binary-search the first true value.

The certified ratios near the end of the new table are numerically consistent
with the OEIS asymptotic comment.  For example,
\[
  \frac{a(67)}{a(66)}\approx 1.947734041055,
\]
and this quotient and \(\exp(2/3)\) both round to \(1.947734041055\) at
twelve decimal places.  This is finite evidence in the certified range, while
the limiting statement itself is supplied by the classical harmonic-number
asymptotic expansion.

\section{Algorithm and Implementation}

The final implementation is the C17 source file \path{src/a226188.c}, built by
\path{src/Makefile}.  It has two supported modes:
\begin{itemize}
\item \texttt{verify-known}, which recomputes and checks the 27 OEIS terms;
\item \texttt{generate}, which performs the known-term replay and then writes
      the certified range through \(n=67\).
\end{itemize}

For each \(n\), the algorithm is:
\begin{enumerate}
\item compute the untrusted locator \(k_0=\lfloor\exp(2n/3-\gamma)\rfloor\);
\item start with radius \(16\), set the trial lower endpoint to
      \(\max(1,k_0-r)-1\), and set the trial upper endpoint to \(k_0+r\);
\item double the radius until the lower endpoint is certified false and the
      upper endpoint certified true for \(H_k>2n/3\), checking unsigned
      addition for overflow at each expansion and exiting without emitting a
      term if such an overflow is detected;
\item binary-search the resulting bracket while \(k_{\rm high}-k_{\rm low}>1\),
      using the overflow-safe midpoint
      \(k_{\rm low}+(k_{\rm high}-k_{\rm low})/2\);
\item certify the final neighbor inequalities
      \(H_{k-1}\leq 2n/3 < H_k\);
\item repeat the final certificate with the higher-order EM6 interval at
      increased precision.
\end{enumerate}

The loop invariant is that the lower endpoint is certified false and the upper
endpoint is certified true.  Since \(H_k\) is strictly increasing, the
invariant proves that the final upper endpoint is the least true value.  The
bracket is finite, and each binary-search step strictly decreases
\(k_{\rm high}-k_{\rm low}\), so the search terminates.  If an interval
comparison is unresolved at 256 bits, the implementation retries at 384 and
then 512 bits; if the comparison is still unresolved, it exits without emitting
a term.  In the completed run no escalation was needed beyond 256 bits.

For \(k\leq 10000\), the program uses directed MPFR summation, the direct-sum
instance of outward interval rounding: lower
accumulators are formed with downward rounding and upper accumulators with
upward rounding, so that the final pair encloses the exact finite sum.  The
threshold \(10000\) is a conservative implementation cutoff: direct summation
is still inexpensive there, while Euler--Maclaurin evaluation is preferable on
the much larger values occurring in the certified range.  For larger values,
the program uses the Euler--Maclaurin interval described above.  The precision
policy is \(256\), \(384\), and \(512\) MPFR bits; the completed run resolved
all terms at 256 bits.  The manifest records MPFR version 4.2.1 and GMP
version 6.3.0.

\section{Certification and Reproducibility}

The production manifest \path{outputs/a226188_manifest.json} reports:
\[
\begin{array}{ll}
\text{range} & 1\leq n\leq 67,\\
\text{known replay} & 1\leq n\leq 27,\ \text{pass},\\
\text{new terms} & 28\leq n\leq 67,\ \text{certified},\\
\text{terms written} & 67,\\
\text{predicate evaluations} & 467,\\
\text{elapsed time} & 0.027904\ \text{s},\\
\text{maximum precision used} & 256\ \text{bits}.
\end{array}
\]
A local final verification run also recorded a wall time of about \(0.03\) s and a
peak resident set size of about \(3184\) KB.  The manifest timing is for the
\texttt{generate} mode, which includes the known-term replay and generation
through \(n=67\).  These timing and memory figures are local telemetry, not a
hardware-normalized benchmark.  The current local verification environment is
an x86-64 GNU/Linux system with an AMD Ryzen 9 7940HS CPU, 64 GB DDR5
installed RAM, GCC 13.3.0, and the Makefile flags
\texttt{-std=c17 -O2} with warning flags
\texttt{-Wall}, \texttt{-Wextra}, \texttt{-Wshadow}, \texttt{-Wconversion},
\texttt{-Wdouble-promotion}, \texttt{-Wformat=2}, and \texttt{-pedantic}.

The validation table \path{outputs/a226188_validation.tsv} contains one
certificate row per index.  Each row records the candidate \(k\), seed,
bracket endpoints, MPFR precision, harmonic method, two comparison outcomes,
positive margins for both neighboring inequalities, checker status, predicate
evaluation count, elapsed time, and final status.  All 67 rows have status
\texttt{CERTIFIED} and checker status \texttt{PASS}.  The smallest recorded
margin for the lower-neighbor inequality is
\[
  2n/3-U_{k-1}=2.633584523432143\cdot 10^{-22}
\]
at \(n=67\), and the smallest recorded margin for the strict crossing
inequality is
\[
  L_k-2n/3=6.228932742447755\cdot 10^{-21}
\]
at \(n=65\).  These margins are taken from the validation TSV and explain why
the 256-bit interval comparisons were resolved in the completed run.

The file \path{data/b226188.txt} is the project's canonical OEIS-ready local
b-file; it contains 67 data rows and the final blank line required by the
local workflow.  Release-level file-integrity records are maintained with the
repository/archive package rather than printed in the paper.

\section{Results}

\begin{theorem}[Certified extension of A226188]
\label{thm:main}
Subject to the recorded source, manifest, and validation table described
above, the values in Table~\ref{tab:newterms} are certified values of A226188
for \(28\leq n\leq 67\).  Together with the reproduced OEIS terms for
\(1\leq n\leq 27\), they give a certified table for \(1\leq n\leq 67\).
For each listed value \(k=a(n)\), the recorded certificate proves
\[
  H_{k-1}\leq 2n/3 < H_k.
\]
\end{theorem}

\begin{proof}[Computational proof]
This is a finite computational certification conditioned on the recorded
source, Makefile, manifest, and validation table.  The public GitHub
repository and Zenodo archive fix the source, build recipe, and released data
artifacts.
The program first recomputes \(a(n)\) for \(1\leq n\leq 27\) and compares the
results with the local OEIS b-file values.  It then processes \(n=28,\ldots,67\)
by certified bracketing and binary search.  Proposition~\ref{prop:certificate}
shows that the final two inequalities for a row are sufficient to identify the
least crossing index.  The validation TSV records both inequalities and their
positive margins for every \(n\), and the manifest records completion of the
run.
\end{proof}

Given a printed pair \((n,k)\), the certificate for that individual row can be
checked independently by evaluating \(H_{k-1}\) and \(H_k\) with any rigorous
arbitrary-precision interval method and verifying the two inequalities in
Proposition~\ref{prop:certificate}.  The source and output artifacts are used
to reproduce the generation and the exact table, not as a substitute for the
mathematical certificate.

\begin{table}[htbp]
\centering
\caption{Certified terms of A226188 for \(28\leq n\leq 67\), extending the 27 previously known terms.}
\label{tab:newterms}
\begin{tabular}{r r r r}
\toprule
\(n\) & \(a(n)\) & \(n\) & \(a(n)\)\\
\midrule
28 & 71804019 & 48 & 44334502845080\\
29 & 139855131 & 49 & 86351820384599\\
30 & 272400600 & 50 & 168190380070122\\
31 & 530563922 & 51 & 327590128640500\\
32 & 1033397411 & 52 & 638058445066582\\
33 & 2012783315 & 53 & 1242768153638596\\
34 & 3920366580 & 54 & 2420581837980561\\
35 & 7635831442 & 55 & 4714649644993433\\
36 & 14872568831 & 56 & 9182883605200051\\
37 & 28967808590 & 57 & 17885814992891026\\
38 & 56421586886 & 58 & 34836810713659947\\
39 & 109894245429 & 59 & 67852842108773716\\
40 & 214044762737 & 60 & 132159290357566703\\
41 & 416902270693 & 61 & 257411148671061652\\
42 & 812014744422 & 62 & 501368456813612867\\
43 & 1581588759549 & 63 & 976532410446924923\\
44 & 3080514265923 & 64 & 1902025418020652442\\
45 & 6000022499693 & 65 & 3704639653630074470\\
46 & 11686448069747 & 66 & 7215652763216299614\\
47 & 22762092724463 & 67 & 14054172515346621401\\
\bottomrule
\end{tabular}
\end{table}

\begin{corollary}
The local b-file \path{data/b226188.txt} extends the A226188 table to 67
indices.  It contains the 27 previously known values followed by the 40 new
certified terms in Table~\ref{tab:newterms}.
\end{corollary}

\section{Domain Limitations}

The computation reported here is bounded by the integer domain of the final
implementation.  The largest certified value is
\[
  a(67)=14054172515346621401,
\]
which fits in \texttt{uint64\_t}.  The manifest records \(n=67\) as the maximum
supported index for this implementation family and records \(n=68\) as
requiring a wider integer type.  This boundary is also visible directly from
the harmonic threshold: an EM4 upper interval, evaluated with the same
outward-rounded interval method used above, gives
\[
  H_{2^{64}-1}<44.938636<\frac{136}{3},
\]
so no \texttt{uint64\_t} value can satisfy the defining inequality for
\(n=68\).  Continuing the sequence should therefore be treated as a new
wider-integer computation, not as a simple extension of the same compiled
artifact.

The paper also does not claim an independent exhaustive reimplementation of the
full non-hit search.  The validation artifacts certify the reported crossing
values by interval arithmetic and replay the 27 known terms.  The finite
extension is reproducible from the stated source and output artifacts, but it
is not a proof of all future terms.

\section{Conclusion}

OEIS A226188 has been extended from 27 known terms to a certified table through
\(n=67\).  The computation confirms every row by the defining inequalities
\[
  H_{a(n)-1}\leq 2n/3 < H_{a(n)}
\]
and records the result in an OEIS-ready local b-file.  The new values are
finite certified data consistent, through the largest \(64\)-bit term, with the
ratio behavior predicted by the harmonic-number asymptotic expansion.  Further
extension begins at a genuine implementation boundary: \(n=68\) requires
integers wider than \texttt{uint64\_t}.

\section*{Data availability}

The code, data files, result tables, manifest, and validation records are part
of the public GitHub repository and Zenodo archive listed in the References.
The essential local artifacts are
\path{src/a226188.c}, \path{src/Makefile}, \path{data/b226188.txt},
\path{outputs/a226188_terms.tsv}, \path{outputs/a226188_validation.tsv},
\path{outputs/a226188_manifest.json}.  The local b-file contains 67 data rows
and ends with the required final blank line.

\section*{AI use disclosure}

Large language models were used as research and drafting assistants during the
A226188 workflow, including OpenAI Codex and external LLM-based review systems
from the ChatGPT, Claude, DeepSeek, Gemini, Grok, Kimi, Qwen, and Z families.
They are named at family level where exact model-version metadata were not
consistently available in the local review artifacts.

Mathematical analysis.  LLM assistance was used to organize the OEIS snapshot,
identify the inverse-harmonic threshold structure, and check that the final
predicate should be \(H_{k-1}\leq 2n/3 < H_k\).  The mathematical acceptance
criterion in this paper is the explicit interval certificate recorded for each
row, together with the author's direct derivation and human-reviewed proof of
Proposition~\ref{prop:certificate}.

Algorithm and code.  LLM assistance was used in the design and review loop for
the C17 implementation using MPFR and GMP.  The final claims rely on the local
source, manifest, validation TSV, known-term replay, and
sanitizer/static-analysis records, not on conversational assertions.

Manuscript.  LLM assistance was used to draft and structure this manuscript
from the project artifacts.  Numerical claims in the manuscript were
checked against the local b-file, result tables, manifest, validation records,
and final verification notes.

The author takes full scientific responsibility for the mathematical
statements, algorithm, implementation, numerical results, validation
boundaries, and conclusions of this paper.

\begin{thebibliography}{9}
\sloppy
\raggedright

\bibitem{oeisA226188}
OEIS Foundation Inc.,
\emph{Sequence A226188: Least positive integer \(k\) such that
\(1+1/2+\cdots+1/k>2n/3\)},
The On-Line Encyclopedia of Integer Sequences,
\url{https://oeis.org/A226188}.
Author-checked OEIS URL; local project snapshot used as source material for
this manuscript.

\bibitem{oeisThresholdFamily}
OEIS Foundation Inc.,
\emph{Related harmonic-threshold sequences A226187, A226189, and A226190},
The On-Line Encyclopedia of Integer Sequences,
\url{https://oeis.org/A226187},
\url{https://oeis.org/A226189},
\url{https://oeis.org/A226190}.
These related entries are cross-referenced in the local A226188 snapshot; the
OEIS URLs were checked by the author.

\bibitem{oeisIntegerThresholds}
OEIS Foundation Inc.,
\emph{Sequences A002387 and A004080: integer-threshold inverse harmonic
numbers},
The On-Line Encyclopedia of Integer Sequences,
\url{https://oeis.org/A002387},
\url{https://oeis.org/A004080}.
These OEIS sources provide the integer-threshold comparison used in the
\(a(3m)=b(2m)\) observation; the OEIS URLs were checked by the author.

\bibitem{boasWrench}
R.~P. Boas, Jr. and J.~W. Wrench, Jr.,
\emph{Partial sums of the harmonic series},
American Mathematical Monthly 78 (1971), no.~8, 864--870.
\url{https://doi.org/10.1080/00029890.1971.11992881}.

\bibitem{dlmf}
NIST Digital Library of Mathematical Functions,
\emph{NIST Digital Library of Mathematical Functions},
\url{https://dlmf.nist.gov/}, Release 1.2.7 of 2026-06-15.
Release date confirmed by the DLMF update note
\url{https://dlmf.nist.gov/about/news/#S1.I1.i1}.
F.~W.~J. Olver, A.~B. Olde Daalhuis, D.~W. Lozier, B.~I. Schneider,
R.~F. Boisvert, C.~W. Clark, B.~R. Miller, B.~V. Saunders, H.~S. Cohl,
and M.~A. McClain, eds.
Section 5.11 was used for the asymptotic expansion and remainder discipline.

\bibitem{mpfr}
L.~Fousse, G.~Hanrot, V.~Lefevre, P.~Pelissier, and P.~Zimmermann,
\emph{MPFR: A multiple-precision binary floating-point library with correct
rounding},
ACM Transactions on Mathematical Software 33 (2007), no.~2, Article 13,
15 pp. \url{https://doi.org/10.1145/1236463.1236468}.

\bibitem{gmp}
GNU MP development team,
\emph{GNU MP: The GNU Multiple Precision Arithmetic Library, version 6.3.0},
\url{https://gmplib.org/manual/}.

\bibitem{projectA226188}
C.~Corti,
\emph{Computational artifacts for OEIS A226188},
GitHub repository \url{https://github.com/carcorti/A226188} and Zenodo archive,
\url{https://doi.org/10.5281/zenodo.20960659}.

\end{thebibliography}

\end{document}
